\section{Implementation}\label{sec:impl}

We implemented $\Pi_{\mathrm{KF}}$, compiled with salted Merkle trees and a Fiat--Shamir transcript, as a self-contained Rust library, \texttt{kronecker-fri}, in the companion repository.
It has about $1{,}650$ lines of library code, including unit tests, and $730$ lines of integration tests, examples and benchmarks (formatted with \texttt{rustfmt}); its only dependency is the BLAKE3 hash function (the parameter check of \cref{rem:pq-params} also uses arbitrary-precision integers, as a development dependency).
The code follows the paper: every function names the statement it implements, and the test suite (\cref{sec:impl:tests}) checks the algebraic statements listed in \cref{app:verification} against direct computations.
This section records the concrete choices; \cref{sec:exp} reports measurements.

\paragraph{Index conventions.}
The paper numbers bits, variables and rounds from $1$; the code numbers them from $0$.
Bit $a_k$ of the paper is bit $k-1$ of the code, the coordinate $z_k$ is \texttt{z[k-1]}, and the challenge $r_j$ is \texttt{rs[j-1]}; the oracle $w_j$, sent in round $j+2$ for $j \in \iv{1}{\ell-1}$, is \texttt{words[j-1]}.
A multilinear polynomial is given by its coefficient form $(\alpha_a)_{a < N}$, which is also the coefficient vector of $U_f$ (\cref{def:kron-enc}), or by its table form, converted with the M\"obius transform.

\subsection{Fields}\label{sec:impl:fields}

\paragraph{Base field.}
Committed polynomials and the commitment word $\Enc(f)$ (whose fibre string is the implicit instance $y$ of \cref{sec:pq}) live over the Goldilocks field $\Fq$, $q = 2^{64} - 2^{32} + 1$ (\cref{ex:goldilocks}).
We take $L = \langle \omega \rangle$ with $\omega = 7^{(q-1)/M}$; since $7$ generates $\Fq^\times$ (the test suite checks $7^{(q-1)/\ell'} \ne 1$ for the six prime factors $\ell'$ of $q - 1 = 2^{32} \cdot 3 \cdot 5 \cdot 17 \cdot 257 \cdot 65537$), $\omega$ has order exactly $M$.
Multiplication reduces a $128$-bit product with $2^{64} \equiv 2^{32} - 1$ and $2^{96} \equiv -1 \pmod q$; elements are kept in canonical form.
This code is shared with the implementation of~\cite{Mkh26}.

\paragraph{Challenge fields.}
The code is generic over the challenge field $\F = \mathbb{F}_{q^e}$ with $e \in \{2, 4\}$.
The quadratic extension is $\Fq[u]/(u^2 - 7)$, used for the classical parameters of \cref{sec:sound:params}; the quartic extension is $\Fq[\iota]/(\iota^4 - 7)$, used for the post-quantum parameters of \cref{rem:pq-params}.
Both polynomials are irreducible because $7$ is not a square in $\Fq$ and $q \equiv 1 \pmod 4$ (\cref{rem:pq-params}); the test suite checks $7^{(q-1)/2} = -1$.
The quartic extension is implemented as the tower $\mathbb{F}_{q^2}[\iota]/(\iota^2 - u)$ with $u = \iota^2$, so that a product costs three products in $\mathbb{F}_{q^2}$ (Karatsuba) and an inverse costs one inverse in $\mathbb{F}_{q^2}$; elements are serialised in the basis $1, \iota, \iota^2, \iota^3$, which is the basis used by the map $\mathrm{num}$ of \cref{sec:pq:ior}.
The test suite checks $\iota^4 = 7$, associativity, distributivity and inverses on random elements, and compares the tower product with the schoolbook product modulo $\iota^4 - 7$.

\subsection{Encoding, opening and folding}\label{sec:impl:alg}

\paragraph{Commit.}
$\Enc(f)$ is computed as in \cref{sec:scheme:enc}: the M\"obius transform if $f$ is given by its table ($\tfrac n2 N$ subtractions, \cref{lem:mobius}), zero-padding to length $M$, and an iterative radix-$2$ NTT over $\Fq$ in natural order (\cref{fact:ntt}).
The NTT is generic over the field of the coefficients, with twiddle factors in $\Fq$, and is also used for $w_A$ over $\F$.

\paragraph{Domain layout.}
Position $i$ of a word on $L_j$ is the point $\omega_j^i$, $\omega_j = \omega^{2^j}$.
Since $\omega_j^{M_{j+1}} = -1$, the fibre over position $k$ of $L_{j+1}$ is the pair of positions $\{k, k + M_{j+1}\}$, and folding reads the two halves of an array and writes one array of half size, with no permutation.
A query index $i \in \iv{0}{M-1}$ is the point $\xi = \omega^i$; at level $j \in \iv{1}{\ell}$ its fold check reads the fibre over position $i \bmod M_j$ of $L_j$.

\paragraph{Opening polynomials.}
The prover computes $G = U_f K_z$ as in the proof of \cref{lem:cost}(3): $n$ successive multiplications by the binomials $X^{2^{k-1}} + z_k$, each a pass of one multiplication and one addition per coefficient.
It reads $v = \coef{N-1}{G}$, $A = (0, \coef{0}{G}, \dots, \coef{N-2}{G})$ and $H = (\coef{N}{G}, \dots, \coef{2N-2}{G}, 0)$ (proof of \cref{lem:split}), and computes $w_A$ by one NTT over $\F$.

\paragraph{Virtual word.}
After $\beta$, the prover computes $\ev_L(K_z)$ by the recursion of \cref{lem:kernel-domain}, and then $h(\xi)$ and $w_0(\xi)$ in one pass over $L$, carrying the powers $\xi$, $\xi^N$ and $\xi^{-(N+1)}$ multiplicatively; $H$ itself is used only for the coefficients of $U_0$.
The verifier computes, at each query point, the powers $\xi^2, \dots, \xi^{2^n} = \xi^N$ ($n$ squarings), $K_{z'}(\xi^2)$ from them, and then $K_z(\pm\xi) = (z_1 \pm \xi)K_{z'}(\xi^2)$ (\cref{lem:kernel-basic}(2)), one inversion $\xi^{-(N+1)} = (\xi^N\cdot\xi)^{-1}$ in $\Fq$, and $h(\pm\xi)$ and $w_0(\pm\xi)$ from the four opened values of $y$ and $w_A$ (\cref{prop:costs}(3)); both use the same function for one value of $h$.

\paragraph{Fold.}
Prover and verifier use the same function for one fibre, in line form (\cref{lem:line}): with $a = w(\xi)$ and $b = w(-\xi)$, $w_e = (a+b)/2$, $w_o = (a-b)(2\xi)^{-1}$ and $\fold_r(w)(\xi^2) = (w_o - w_e) + r\,(2w_e - w_o)$; this is the count of $7$ operations of \cref{prop:costs}.
The prover reads $(2\xi)^{-1}$ from one table of $\omega^{-t}$, $t < M/2$, shared by all levels.
The verifier obtains $(2\omega_{j-1}^{k})^{-1} = \tfrac12\,\omega^{M - k 2^{j-1}}$ by one exponentiation in $\Fq$, without inversion.
The final polynomial is $P = \pfold^{(\ell)}_r(U_0)$, computed on coefficients (\cref{lem:iterated}), and the verifier evaluates it by Horner's rule.

\subsection{Commitments and transcript}\label{sec:impl:merkle}

\paragraph{Salted Merkle trees.}
Leaves are hashed as $\mathsf{H}(\mathtt{0x00} \,\|\, \text{data} \,\|\, \gamma)$ with a salt $\gamma$ of $32$ bytes, and internal nodes as $\mathsf{H}(\mathtt{0x01} \,\|\, \text{left} \,\|\, \text{right})$, with $\mathsf{H} = $ BLAKE3 and $\lambda_{\mathsf{H}} = 256$.
The salts of a tree are derived from a $32$-byte seed supplied by the caller, by BLAKE3 in keyed mode; the seed must be fresh and secret for every commitment and every proof, since the salts are what the compiler of~\cite{CDHZ25} requires to be random.
Setting the salt length to $0$ gives unsalted trees, which we use only to measure the cost of salting (\cref{sec:exp}).

\paragraph{Fibre leaves.}
As in \cref{rem:layout}, every tree has one leaf per fibre: leaf $k$ of the tree of a word $u$ on $L_j$ holds $(u(\omega_j^k), u(-\omega_j^k))$, $k < M_j/2$, that is, the symbol $\mathrm{fib}_j(u)$ at position $\omega_{j+1}^k$ with the representative $\xi_\eta = \omega_j^k$ (\cref{sec:pq:ior}).
Leaves of the commitment tree hold two elements of $\Fq$ ($16$ bytes), and the other leaves two elements of $\F$.
Each query point opens one leaf in the commitment tree, one in the tree of $w_A$, and one in each tree of $w_1, \dots, w_{\ell-1}$, with its salt and authentication path; the value $w_j(\xi^{2^j})$ compared at level $j$ is one of the two values of the leaf opened at level $j+1$.
The library also supports openings of aligned groups of $2^g$ leaves with one authentication path; the protocol uses $g = 0$.

\paragraph{Fiat--Shamir.}
The transcript is a BLAKE3 hash chain. It absorbs the parameters $(n, R, \ell, \kappa)$, the salt length and $e$, the commitment root, $z$ and $v$.
In round $1$ it absorbs the root of $w_A$ and a round salt, and derives $\beta$; in round $2$ it absorbs a round salt and derives $r_1$; in round $j+1$, $j \in \iv{2}{\ell}$, it absorbs the root of $w_{j-1}$ and a round salt, and derives $r_j$; in round $\ell+2$ it absorbs the coefficients of $P$ and a round salt, and derives the query indices.
A challenge in $\F$ is the image under the map $\phi$ of \cref{sec:pq:ior} of $\lambda_{\mathrm{ch}} = 64(e+1)$ output bits ($192$ for $e = 2$, $320$ for $e = 4$): the bits are read as an integer, whose first $e$ digits in base $q$ are the coordinates.
The $\kappa$ query indices are $\kappa$ consecutive blocks of $n+R$ output bits (the map $\mathrm{pos}$).

\paragraph{Relation to the compiler of~\cite{CDHZ25}.}
The implementation has the structure of $\mathsf{BCS}[\Pi^\Sigma_{\mathrm{KF}}]$: salted trees, one round salt per round, challenges derived from the roots and salts of all previous rounds, and the maps $\phi$ and $\mathrm{pos}$.
It departs from \cref{def:pikf-sigma} in one inessential point: round $2$ sends no string, so no tree is built for the dummy symbol.
We have not checked the byte-level encoding of the transcript and of the salts against~\cite[Construction~11.7]{CDHZ25}, and we therefore do not claim that \cref{cor:pq} applies to the code as it stands (\cref{rem:scope}).

\paragraph{Proof format.}
A proof consists of the root of $w_A$, the roots of $w_1, \dots, w_{\ell-1}$, the $\ell + 2$ round salts, the $N_\ell$ coefficients of $P$, and, per query point, the openings described above: $2$ elements of $\Fq$, $2\ell$ elements of $\F$, $\ell+1$ leaf salts, and $\ell + 1$ authentication paths, two of length $\log_2 M - 1$ and one of length $\log_2 M_j - 1$ for each $j \in \iv{1}{\ell-1}$.
\Cref{sec:exp} reports the measured sizes.

\subsection{Testing}\label{sec:impl:tests}

The test suite checks each component against a direct definition, and the protocol end to end.
\begin{itemize}
  \item \emph{Arithmetic.} Base-field operations against $128$-bit integer arithmetic; the extensions as above; the NTT against Horner evaluation, over $\mathbb{F}_{q^4}$.
  \item \emph{\Cref{sec:kernel}.} For $n \le 8$ and both extensions: the coefficients of $K_z$ against the product of binomials of \cref{def:kernel} (\cref{lem:kernel-basic}(3)); $K_z(\xi)$ against Horner evaluation; the kernel product against schoolbook multiplication; the extraction identity (\cref{thm:extraction}) against $f(z) = \sum_a \alpha_a z^a$; the M\"obius transform against the table; and the split identity \eqref{eq:identity} with $A(0) = 0$ and $\coef{N-1}{H} = 0$ at a random point.
  \item \emph{\Cref{sec:fold}.} For $n \le 7$ and $\ell = n$: the explicit formula \eqref{eq:fold-explicit} against the line form, the consistency of word and polynomial folds (\cref{prop:fold-consistency}), and the coefficient formula of \cref{lem:iterated}, after every fold.
  \item \emph{Transcript and trees.} Group openings of every size and position, with and without salts, and rejection of a modified leaf; canonical challenge coordinates.
  \item \emph{Completeness.} Honest proofs verify for all $1 \le n \le 9$ over $\mathbb{F}_{q^2}$ and $1 \le n \le 7$ over $\mathbb{F}_{q^4}$, all $1 \le \ell \le n$, rates $\rho \in \{1/4, 1/8, 1/16\}$, with and without salts; and the returned value equals $f(z)$ computed from the table form by $\sum_b f(b) \prod_k (b_k z_k + (1-b_k)(1-z_k))$, independently of the M\"obius transform and of the kernel product.
  \item \emph{Virtual word.} For $n \le 6$ and all $\xi \in L$: the virtual word of honest openings equals $H(\xi)$ (\cref{lem:virtual}(3)), and $\ev_L(K_z)$ computed by \cref{lem:kernel-domain} equals $K_z(\xi)$.
  \item \emph{Tampering.} Proofs are rejected after changing the claimed value, the point, the commitment, an opened value of $y$, $w_A$ or $w_1$, a coefficient of $P$, a root, a round salt, a leaf salt or an authentication path, and when a proof for another polynomial is checked against the original commitment.
  \item \emph{Two cheating provers.} Both commit honestly and claim $v + 1$, and both are rejected by a fold check in every tested configuration. The first sends the honest $w_A$; its virtual word is $\ev_L(H) - \ev_L(X^{-1})$, far from $\Cc_0$ (\cref{rem:design}(1)). The second sends $w_A$ for $A - X^N$, so that its virtual word is exactly $\ev_L(H)$; now $w_A$ has degree $N$ and the batched word is far from $\Cc_0$. This is the situation of \cref{rem:why-x}: the degree bound $N$ of the proximity test is what prevents $A$ from absorbing the error in $v$.
\end{itemize}
Independently, \texttt{examples/small\_field\_checks.rs} checks the extraction identity, the split lemma, the identity \eqref{eq:identity} at every point of $L$ for honest openings together with the folds, and the root bound $2N$ of \cref{lem:identity} for a wrong value, over $\mathbb{F}_{257}$, with $1 \le n \le 5$ and $20$ random instances each; and \texttt{examples/pq\_params.rs} evaluates the bound of \cref{rem:pq-params} exactly.
